% !TEX program = xelatex
% 공통수학2 · Ⅲ. 함수와 그래프 · 중단원(2.유리함수와 무리함수)·대단원(Ⅲ) 마무리 · 기말 총정리
\documentclass[11pt,a4paper]{article}
\usepackage{kotex}
\usepackage{iftex}
\ifPDFTeX\else
  \setmainhangulfont{Noto Serif CJK KR}
  \setsanshangulfont{Noto Sans CJK KR}
\fi
\usepackage[margin=2.1cm]{geometry}
\usepackage{parskip}
\usepackage{enumitem}
\usepackage{array}
\usepackage{tabularx}
\usepackage{booktabs}
\usepackage{xcolor}
\usepackage{amssymb}
\usepackage{tikz}
\usepackage[most]{tcolorbox}
\usepackage{fancyhdr}
\usepackage{titlesec}

\definecolor{goldc}{HTML}{B07C22}
\definecolor{goldstrong}{HTML}{8A5F16}
\definecolor{indigoc}{HTML}{2B3B57}
\definecolor{indigosoft}{HTML}{6E82A6}
\definecolor{linec}{HTML}{D6DACB}
\definecolor{surface2}{HTML}{E4E8DC}
\definecolor{notebg}{HTML}{FBF3E3}
\definecolor{inksoft}{HTML}{52604F}

\titleformat{\section}{\normalfont\sffamily\bfseries\large\color{indigoc}}{\thesection}{0.6em}{}
\titlespacing{\section}{0pt}{16pt}{8pt}
\newtcolorbox{ideabox}{enhanced, breakable, colback=white, colframe=goldc, boxrule=0.5pt, leftrule=3pt, arc=2pt, left=10pt, right=10pt, top=8pt, bottom=8pt}
\newtcolorbox{calloutbox}{enhanced, breakable, colback=white, colframe=indigosoft, boxrule=0.5pt, leftrule=3pt, arc=2pt, left=10pt, right=10pt, top=7pt, bottom=7pt}
\newtcolorbox{answerbox}{enhanced, breakable, colback=white, colframe=linec, boxrule=0.5pt, arc=2pt, left=10pt, right=10pt, top=7pt, bottom=7pt}
\newcommand{\lessonbanner}[3]{%
  \begin{tcolorbox}[enhanced, colback=indigoc, colframe=indigoc, boxrule=0pt, arc=0pt,
    left=14pt, right=14pt, top=14pt, bottom=10pt, borderline south={2.4pt}{0pt}{goldc}, fontupper=\color{white}]
    {\sffamily\footnotesize\color{white!85!black} #1}\\[4pt]{\sffamily\bfseries\Large #2}\\[3pt]{\sffamily\small\color{white!88!black} #3}
  \end{tcolorbox}}
\pagestyle{fancy}\fancyhf{}\renewcommand{\headrulewidth}{0pt}
\fancyfoot[L]{\footnotesize\color{inksoft} 공통수학2 · 2026학년도 2학기 · 지평선고등학교}
\fancyfoot[R]{\footnotesize\color{inksoft} 교사용 지도안 -- 배포 금지}

\begin{document}
\lessonbanner{공통수학2 · Ⅲ. 함수와 그래프}{중단원(2. 유리함수와 무리함수)\textperiodcentered 대단원(Ⅲ) 마무리 \textperiodcentered\ 기말 총정리 --- 42차시}{교사용 지도안 (2026학년도 2학기)}
\vspace{10pt}

\noindent\begin{tabularx}{\linewidth}{@{}>{\bfseries\color{indigoc}}p{2.6cm}X@{}}
\toprule
대상 & 고1 · 2개 반 · 반당 13명 \\
차시 & 50분 (도입 10 · 전개 30 · 정리 10) \\
성취기준 & 함수, 합성함수, 역함수, 유리함수, 무리함수를 통합적으로 이해하고 적용할 수 있다. \\
선행 학습 & 33$\sim$41차시 --- 대단원 Ⅲ `함수와 그래프' 전체 \\
\bottomrule
\end{tabularx}

\section{학습 목표 \& Big Idea}
\begin{ideabox}
\textbf{\color{goldstrong}영속적 이해 (Big Idea)}\\
대단원 Ⅲ 전체를 관통하는 하나의 아이디어는, 복잡해 보이는 함수의 그래프도 결국 \textbf{가장 단순한 원형을 평행이동한 것}이라는 점이다. 유리함수는 $y=\dfrac kx$의, 무리함수는 $y=\sqrt{ax}$의 평행이동이며, 이는 이미 한 학기 동안 배운 원의 방정식·평행이동의 원리와 같은 논리다.
\vspace{4pt}
\small\color{inksoft} 핵심개념: \textbf{대응} \quad\textbullet\quad 핵심개념: \textbf{원형과 평행이동} \quad\textbullet\quad 일반화: \textbf{복잡한 그래프도 단순한 원형의 변형일 뿐이다}
\end{ideabox}
\begin{calloutbox}
\textbf{차시 학습 목표}
\begin{itemize}[leftmargin=1.4em, itemsep=2pt, topsep=4pt]
  \item 함수의 뜻, 합성함수, 역함수, 유리함수, 무리함수를 종합 문제를 통해 점검하고 정리할 수 있다.
  \item 한 학기 동안 배운 대단원 Ⅰ`도형의 방정식', Ⅱ`집합과 명제', Ⅲ`함수와 그래프'가 어떻게 서로 연결되는지 설명할 수 있다.
\end{itemize}
\end{calloutbox}

\section{질문의 연결}
\begin{tabularx}{\linewidth}{@{}>{\bfseries\color{indigoc}\small}p{2.4cm}X@{}}
사실적 질문 & 대단원 Ⅲ에서 배운 하위 개념들을 각각 한 문장으로 요약하면 무엇일까? \\[6pt]
개념적 질문 & `평행이동'이라는 하나의 아이디어가 원의 방정식, 유리함수, 무리함수에서 어떻게 반복해서 나타날까? \\[6pt]
논쟁적 질문 & 한 학기 동안 배운 개념 중 앞으로 실생활이나 다른 과목에서 가장 자주 마주치게 될 것은 무엇이라고 생각하는가? \\
\end{tabularx}

\section{수업 흐름}
\begin{tabularx}{\linewidth}{@{}>{\centering\arraybackslash}p{1.6cm}X>{\raggedright\arraybackslash\small}p{3.1cm}@{}}
\toprule
\textbf{단계} & \textbf{활동 \& 발문} & \textbf{ATL / VTR} \\
\midrule
\textcolor{goldstrong}{\textbf{도입}}\newline\footnotesize 10분 &
\textbf{마음공부 (2분)} --- 한 학기를 마무리하는 유무념 대조.\newline
\textbf{개념 지도 그리기 (8분)} --- 대단원 Ⅲ의 하위 개념(함수, 여러 가지 함수, 합성함수, 역함수, 유리함수, 무리함수)을 화이트보드에 자유롭게 배치하며 `평행이동'이라는 공통 실마리로 화살표 연결. &
ATL: 사고(구조화)\newline VTR: Connect-Extend-Challenge \\
\addlinespace
\textcolor{goldstrong}{\textbf{전개}}\newline\footnotesize 30분 &
\textbf{종합 문제 풀이 (30분)} --- 문제 1(함수의 뜻, 정의역·공역·치역) $\to$ 문제 2(합성함수와 역함수) $\to$ 문제 3(유리함수) $\to$ 문제 4(무리함수)를 모둠별로 순환하며 해결. 각 문제가 어느 차시 개념과 연결되는지 짚어 준다. &
ATT: 촉진자 역할\newline ATL: 협업(모둠 순환)\newline VTR: 개념 연결 확인 \\
\addlinespace
\textcolor{goldstrong}{\textbf{정리}}\newline\footnotesize 10분 &
\textbf{개념 연결 질문 (5분)} --- ``대단원 Ⅲ에서 가장 헷갈렸던 개념은 무엇이었고, 오늘 다시 보니 어떻게 정리되었나?''\newline
\textbf{성찰 (5분)} --- 감사 일기 1문장 + 대단원 Ⅲ `함수와 그래프' 총마무리 + 한 학기 전체(Ⅰ$\sim$Ⅲ) 회고 + 기말고사 응원. &
마음공부 · 감사 일기 \\
\bottomrule
\end{tabularx}

\begin{calloutbox}
\textbf{지도상의 유의점} --- 이 차시는 새로운 개념을 배우는 시간이 아니라 한 학기 전체를 아우르는 \textbf{연결}을 확인하는 마지막 시간임을 분명히 한다. 아래 `학기 전체 핵심 공식 총정리' 표를 활용해 기말고사 범위(Ⅰ$\sim$Ⅲ 전체)를 학생 스스로 점검하도록 안내한다.
\end{calloutbox}

\section{형성평가 \& 답안}
\begin{minipage}[t]{0.48\linewidth}
\begin{answerbox}
\textbf{문제 1} \quad $X=\{1,2,3\}$, $Y=\{1,2,\ldots,9\}$, $f(x)=x^2$일 때, 정의역·공역·치역을 구하고 일대일함수인지 판별하시오.\\[4pt]
\textbf{\color{goldstrong}답} \; 정의역 $=\{1,2,3\}$, 공역 $=\{1,\ldots,9\}$, 치역 $=\{1,4,9\}$. 서로 다른 원소의 함숫값이 모두 다르므로 일대일함수.
\end{answerbox}
\end{minipage}\hfill
\begin{minipage}[t]{0.48\linewidth}
\begin{answerbox}
\textbf{문제 2} \quad $f(x)=2x-1$, $g(x)=x+2$일 때, $(f\circ g)(x)$와 $(f\circ g)^{-1}(x)$를 각각 구하시오.\\[4pt]
\textbf{\color{goldstrong}답} \; $(f\circ g)(x)=2x+3$. \; $(f\circ g)^{-1}(x)=\dfrac{x-3}{2}$
\end{answerbox}
\end{minipage}

\vspace{6pt}
\begin{minipage}[t]{0.48\linewidth}
\begin{answerbox}
\textbf{문제 3} \quad $y=\dfrac{3x-1}{x-2}$를 $y=\dfrac{k}{x-p}+q$ 꼴로 변형하고, 점근선과 정의역을 구하시오.\\[4pt]
\textbf{\color{goldstrong}답} \; $3x-1=3(x-2)+5$이므로 $y=\dfrac{5}{x-2}+3$. 점근선 $x=2$, $y=3$. 정의역 $\{x\mid x\neq2\}$
\end{answerbox}
\end{minipage}\hfill
\begin{minipage}[t]{0.48\linewidth}
\begin{answerbox}
\textbf{문제 4} \quad $y=\sqrt{3-x}+2$의 정의역과 치역을 구하시오.\\[4pt]
\textbf{\color{goldstrong}답} \; $3-x\geq0 \Rightarrow x\leq3$, 정의역 $\{x\mid x\leq3\}$. 치역 $\{y\mid y\geq2\}$
\end{answerbox}
\end{minipage}

\section{학기 전체 핵심 공식 총정리 (기말고사 대비)}
\begin{tabularx}{\linewidth}{@{}>{\bfseries\small}p{3.6cm}X@{}}
\toprule
\textbf{대단원} & \textbf{핵심 공식·개념} \\
\midrule
Ⅰ. 도형의 방정식 & 내분점 공식, 두 직선의 평행·수직 조건, 점과 직선 사이의 거리, 원의 방정식(표준형·일반형), 원과 직선의 위치 관계(판별식), 접선의 방정식, 평행이동 $f(x-a,y-b)=0$, 대칭이동(축·원점·$y=x$) \\
Ⅱ. 집합과 명제 & 집합의 연산(합·교·여·차, 드모르간 법칙), 명제와 조건, 진리집합, 역과 대우, 충분·필요조건, 대우증명법·귀류법, 절대부등식 \\
Ⅲ. 함수와 그래프 & 함수의 뜻과 그래프, 일대일함수·일대일대응, 합성함수, 역함수와 그 그래프($y=x$ 대칭), 유리함수($y=\dfrac{k}{x-p}+q$), 무리함수($y=\sqrt{a(x-p)}+q$) \\
\bottomrule
\end{tabularx}

\section{성취수준 (대단원 Ⅲ 종합)}
\begin{tabularx}{\linewidth}{@{}>{\bfseries\color{goldstrong}}p{0.9cm}X@{}}
\toprule
A & 함수와 그래프의 모든 개념을 통합적으로 이해하고, 근거를 들어 스스로 설명하며 다양한 문제에 적용할 수 있다. \\
B & 함수와 그래프의 주요 개념을 이해하고 대부분의 문제를 스스로 해결할 수 있다. \\
C & 함수와 그래프의 기본 개념을 알고, 안내에 따라 문제를 해결할 수 있다. \\
D & 함수와 그래프의 일부 개념을 알고, 안내된 절차에 따라 간단한 문제를 해결할 수 있다. \\
E & 함수와 그래프의 기본 용어를 안다. \\
\bottomrule
\end{tabularx}

\section{다음 차시}
\begin{calloutbox}
\textbf{기말고사.} 대단원 Ⅲ `함수와 그래프'가 여기서 마무리되며, 이는 곧 2026학년도 2학기 공통수학2 전체(대단원 Ⅰ$\sim$Ⅲ)의 마지막 수업이다. 위 `학기 전체 핵심 공식 총정리' 표를 중심으로 기말고사를 준비하도록 안내한다.
\end{calloutbox}
\end{document}
